\begingroup
\begin{table}[H]
    \centering
    \small
    \caption{
    Few-shot exemplars for full chain of thought prompt for math word problems. These exemplars are the same as in \cref{tab:appendix-mwp-prompts}, except that the chains of thought were written by a different annotator (``Annotator B'' instead of ``Annotator A'').
    Annotators were co-authors and familiar with the goal of chain of thought prompting.
    }
    \begin{tabular}{p{\linewidth}}
        \toprule
        \underline{\textbf{\textsc{Prompt for Math Word Problems}}} \\
        \vspace{-2mm}
        \textbf{Q:} There are 15 trees in the grove. Grove workers will plant trees in the grove today. After they are done, there will be 21 trees. How many trees did the grove workers plant today? \\
        \vspace{-1mm}
        \textbf{A:} \hl{There are 21 trees now and there are 15 trees in the beginning, so the workers plant 21 - 15 = 6 trees.} The answer is 6. \\
        \vspace{0mm}
        \textbf{Q:} If there are 3 cars in the parking lot and 2 more cars arrive, how many cars are in the parking lot? \\
        \vspace{-1mm}
        \textbf{A:} \hl{ There are 3 cars in the beginning, 2 more arrive, so now there should be 3 + 2 = 5 cars.}  The answer is 5. \\
        \vspace{0mm}
        \textbf{Q:} Leah had 32 chocolates and her sister had 42. If they ate 35, how many pieces do they have left in total? \\
        \vspace{-1mm}
        \textbf{A:} \hl{Leah had 32 chocolates and her sister had 42, in total they have 32 + 42 = 74 chocolates. After they ate 35, now there are 74 - 35 = 39 chocolates.} The answer is 39. \\
        \vspace{0mm}
        \textbf{Q:} Jason had 20 lollipops. He gave Denny some lollipops. Now Jason has 12 lollipops. How many lollipops did Jason give to Denny? \\
        \vspace{-1mm}
        \textbf{A:} \hl{Jason started with 20 lollipops, but now he only has 12, so he gave Denny 20 - 12 = 8 lollipops.} The answer is 8. \\
        \vspace{0mm}
        \textbf{Q:} Shawn has five toys. For Christmas, he got two toys each from his mom and dad. How many toys does he have now? \\
        \vspace{-1mm}
        \textbf{A:} \hl{Shawn got 2 toys each from his mom and dad, so he got 2 * 2 = 4 more, now he will have 5 + 4 = 9 toys.} The answer is 9. \\
        \vspace{0mm}
        \textbf{Q:} There were nine computers in the server room. Five more computers were installed each day, from monday to thursday. How many computers are now in the server room? \\
        \vspace{-1mm}
        \textbf{A:} \hl{5 computers were installed from monday to thursday, so in total 5 * 4 = 20 computers are installed. 9 computers are there in the beginning, so now there are 20 + 9 = 29 computers.} The answer is 29. \\
        \vspace{0mm}
        \textbf{Q:} Michael had 58 golf balls. On tuesday, he lost 23 golf balls. On wednesday, he lost 2 more. How many golf balls did he have at the end of wednesday? \\
        \vspace{-1mm}
        \textbf{A:} \hl{Michael started with 58 golf balls and lost 23, so he has 58 - 23 = 35. After he lost 2 more, he has 35 - 2 = 33 balls now.} The answer is 33. \\
        \vspace{0mm}
        \textbf{Q:} Olivia has \$23. She bought five bagels for \$3 each. How much money does she have left? \\
        \vspace{-1mm}
        \textbf{A:} \hl{5 bagels for \$3 each should cost 5 * 3 = 15 dollars. Olivia had \$23 in the beginning, so now she has 23 - 15 = 8 dollars left.} The answer is 8. \\
        \bottomrule
    \end{tabular}
    \label{tab:appendix-mwp-prompts-alt}
\end{table}
\endgroup